
Dataset characteristic extraction from published benchmarks proves more challenging than anticipated. This work quantifies a 14\% average metadata coverage barrier that prevents meta-learning approaches from being tested fairly on literature-derived training data.

Machine learning practitioners selecting methods for new datasets face overwhelming algorithmic choices. Should one apply linear models, polynomial basis expansions, recurrent architectures, or data augmentation? Current practice relies on trial-and-error or domain folklore ($vision\rightarrow CNN$, time-$series\rightarrow RNN)$, neither providing systematic evidence-based guidance. Zhou et al. report that no single algorithm achieves optimal performance across nine medical federated learning datasets, while Champneys et al. demonstrate NARX-Poly outperforming LSTM on structured nonlinear system identification (0.032 vs. 0.126 RMSE on Wing-Hing saturation). These observations suggest systematic dataset-method relationships exist but remain unformalized.

This variability motivates a meta-learning hypothesis: training a classifier on aggregated benchmark results (50-60 datasets) using fast-to-compute characteristics would enable prediction of top-performing method families for new datasets. The approach requires a causal chain: (1) dataset characteristics correlate with method family performance, (2) published benchmarks provide sufficient training examples, (3) a Random Forest meta-classifier learns generalizable patterns, and (4) predicted methods achieve competitive performance on held-out datasets.

Experiments revealed a practical bottleneck that prevented hypothesis testing: literature mining provides insufficient metadata for meta-learning. While 29 benchmarks were collected from accessible sources (OGB, FedML, LEAF, manual paper extraction), verifying that data sources exist, feature extraction from APIs and documentation yielded only 13.8\% coverage for sample size and 0\% for dimensionality. This sparsity cascaded through subsequent experiments: correlation analysis (h-m1) found zero significant feature-method relationships due to insufficient feature diversity, and meta-classifier training (h-m2) achieved 25.6\% accuracy with a degenerate feature set containing one usable dimension after NaN filtering, performing worse than a 48.3\% majority-class baseline.

The negative results do not refute meta-learning for method selection. All three sub-hypothesis failures were classified as data limitations rather than hypothesis invalidations. Data collection (h-e1) achieved proof-of-concept validation by confirming source accessibility (OGB library loads, GitHub READMEs fetch, paper tables exist). Correlation analysis (h-m1) could not compute relationships because only four sample size values were extracted, not because correlations were tested and found insignificant. Meta-classifier training (h-m2) failed because only one feature had sufficient coverage (num\_classes: 75.9\%), creating degenerate input that no learning algorithm could exploit.

The findings expose an infrastructure gap: meta-learning requires two-stage data collection. Stage 1 (identify benchmark sources) succeeded, confirming that OGB, FedML, LEAF, and published papers are accessible. Stage 2 (extract dataset characteristics) is the bottleneck---APIs provide identifiers but not feature statistics, READMEs document usage but not sample sizes, and paper tables report method rankings but omit dataset properties. Extracting characteristics requires downloading raw datasets and computing statistics, an engineering effort beyond literature mining.

This work makes three contributions. First, it quantifies the metadata bottleneck: literature mining alone yields 14\% average coverage across critical features (sample size, dimensionality, class balance), insufficient for correlation analysis or meta-learning training. Second, it demonstrates transparent negative result reporting: all deviations were classified as data limitations or scope reductions, verified by mock-data-elimination protocols and coverage analysis. Third, it proposes two-stage collection as a field-wide infrastructure need, reframing the research question from testing whether meta-learning works to identifying what data infrastructure must exist for fair evaluation.

The paper is organized as follows. Section 2 reviews meta-learning, benchmark studies, and method selection heuristics. Section 3 describes the three-stage experimental design (data collection $\rightarrow$ correlation $\rightarrow$ meta-classifier training) and rationale. Section 4 details experimental protocols for h-e1, h-m1, and h-m2. Section 5 presents results showing the metadata bottleneck and cascading effects. Section 6 discusses implications and proposes community-driven benchmark metadata repositories. Section 7 concludes with a call for infrastructure investment to unlock meta-learning research.
